%% LyX 2.4.4 created this file.  For more info, see https://www.lyx.org/.
%% Do not edit unless you really know what you are doing.
\documentclass[english]{article}
\usepackage[T1]{fontenc}
\usepackage[latin9]{inputenc}
\usepackage{units}
\usepackage{enumitem}
\usepackage{amsmath}
\usepackage{amsthm}
\usepackage{graphicx}
\usepackage{geometry}
\geometry{verbose,tmargin=1in,bmargin=1in,lmargin=1in,rmargin=1in}

\makeatletter
%%%%%%%%%%%%%%%%%%%%%%%%%%%%%% Textclass specific LaTeX commands.
\numberwithin{equation}{section}
\numberwithin{figure}{section}
\newlength{\lyxlabelwidth}      % auxiliary length 

%%%%%%%%%%%%%%%%%%%%%%%%%%%%%% User specified LaTeX commands.
\usepackage{fancyhdr}
\usepackage{lastpage}
\pagestyle{fancy}
\usepackage{enumitem}
\usepackage{ifthen}
%sets math fonts to be more readable
\everymath=\expandafter{\the\everymath\displaystyle}
%\DeclareMathSizes{10}{10}{10}{10}
\usepackage{multicol}
% Added by lyx2lyx
\setlength{\parskip}{\smallskipamount}
\setlength{\parindent}{0pt}

\makeatother

\usepackage{babel}
\begin{document}
\renewcommand{\author}{Dr.\,James Ashe}

\newcommand{\classNum}{232}

\newcommand{\classSec}{C}

\newcommand{\nameType}{Name:}

\newcommand{\examType}{Final Exam}

\renewcommand{\date}{January $16^{\text{th}}$ 10am-12pm, 2012} %%\today

%%First page's header%%%%%%%%%%%%%%%%%%%%%%
\fancypagestyle{firststyle} %%header settings for first pagr
{
	\fancyhead[L]{MTH \classNum{}(\classSec{}) \author{}\\
	\textbf{\examType{}} \date{}}
	\fancyhead[C]{\textbf{\nameType}}
	\setlength{\headsep}{14pt}
}
%%%%%%%%%%%%%%%%%%%%%%%%%%%%%%%%%%%%%%%%%%

%%Next page's header%%%%%%%%%%%%%%%%%%%%%%%
\setlist{nolistsep}
\lhead{MTH \classNum{}(\classSec{})} 
\chead{\textbf{\nameType}} 
\cfoot{\thepage{}~of~\pageref{LastPage}} 
\setlength{\headsep}{5pt} 
%%%%%%%%%%%%%%%%%%%%%%%%%%%%%%%%%%%%%%%%%%%

%%Various settings; see the preamble for other settings%%%%%
%%\setenumerate[0]{leftmargin=12pt,itemindent=0pt,labelwidth=0pt}
\renewcommand{\labelenumi}{\textbf{\arabic{enumi}.}}
\renewcommand{\labelenumii}{\textbf{\alph{enumii}.}}
\thispagestyle{firststyle}
%%%%%%%%%%%%%%%%%%%%%%%%%%%%%%%%%%%%%%%%%%

\emph{Complete the following steps for $f(x)=x^{2}+1$, interval $[0,2]$,
and $n=4$.}
\begin{enumerate}
\item Sketch the graph of the function on the given interval and illustrate
the \emph{left Riemann sums }using $n$ subintervals.\vspace{4cm}
\item Write an expression that gives the approximate value of the left Riemann
if $n$ is very large (you do not have to compute the value).\vfill 
\end{enumerate}
\emph{Suppose that ${\displaystyle \int_{0}^{3}f(x)\,dx=2}$ and ${\displaystyle \int_{3}^{6}f(x)\,dx=-5}$,
and ${\displaystyle \int_{3}^{6}g(x)\,dx}=1$. Evaluate the following
integrals.}
\begin{enumerate}[resume]
\item ${\displaystyle \int_{3}^{0}f}(x)\,dx$\vfill
\item ${\displaystyle \int_{3}^{6}\left(2f(x)-g(x)\right)\,dx}$\vfill
\item ${\displaystyle \int_{0}^{6}f(x)\,dx}$\vfill
\end{enumerate}
\emph{Write an expression which gives the average value of $f(z)$
on the given interval (you do not need to compute the value)}.\emph{ }
\begin{enumerate}[resume]
\item \emph{$f(z)=\sin(2z)e^{1-\cos(2\pi)}$} on the interval $[-\pi,\pi]$
\vfill
\end{enumerate}
\newpage

\emph{The velocity of an object is, $\frac{ds}{dt}=-20t$ with intial
position $s(0)=100$. Use this information to solve the following
two problems.}
\begin{enumerate}[resume]
\item Find the position function $s(t)$.\vfill
\item Find the displacement of the object from $t=1$ to $t=2$.\vspace{3cm}
\end{enumerate}
\emph{Evaluate the area of the region described.}
\begin{enumerate}[resume]
\item Write an integral expression which gives the area of the shaded region
in the figure.\\
\includegraphics[width=3cm]{images/2(8-x).svg}
\item Find the area of the region bounded by the graphs of $f(x)=x^{2}+2$
and $g(x)=x-2$ from $x=-1$ to $x=1$.\vfill
\end{enumerate}
\emph{Let $R$ be the region bounded by the following curves. Use
the disk method to find the volume of the solid generated when $R$
is revolved about the $x$-axis.}
\begin{enumerate}[resume]
\item $y=\sqrt{x+1}$ from $x=0$ to $x=3$.\vfill
\end{enumerate}
\newpage

\emph{Find the length of the curve for the given function and interval.}
\begin{enumerate}[resume]
\item $y=\frac{2}{3}x^{\nicefrac{3}{2}}+9$ from $x=0$ to $x=3$.\vfill
\end{enumerate}

\subsubsection*{\textmd{\emph{Evaluate the following integrals. }}}
\begin{enumerate}[resume]
\item ${\displaystyle \int}_{-1}^{1}6x^{5}-x+5\,dx$\vfill
\item ${\displaystyle \int20(5x+1)^{-4}\,dx}$\vfill
\item ${\displaystyle \int\frac{t^{5}-5t+2}{t^{2}}\,dt}$\vfill
\end{enumerate}
\newpage
\begin{enumerate}[resume]
\item ${\displaystyle \int\left(\frac{2}{x}-4e^{2x}\right)\,dx}$\vfill
\end{enumerate}
\emph{Evaluate the following integrals by using integration by parts.}
\begin{enumerate}[resume]
\item ${\displaystyle \int}xe^{-x}\,dx$\vfill
\item ${\displaystyle \int5x\sin x\,dx}$\vfill
\end{enumerate}
\emph{Evaulate the integrals.}
\begin{enumerate}[resume]
\item ${\displaystyle \int\cos^{2}x\,dx}$\vfill
\end{enumerate}
\newpage
\begin{enumerate}[resume]
\item ${\displaystyle \int\sin^{3}x}\cos^{2}x\,dx$\vfill
\item ${\displaystyle \int\frac{6}{x^{2}+2x-8}\,dx}$\vfill
\item ${\displaystyle \int\frac{dx}{\sqrt{9-x^{2}}}}$\vspace{1cm}\vfill
\end{enumerate}
\newpage
\begin{enumerate}[resume]
\item ${\displaystyle \int_{0}^{2}\sqrt{4-x^{2}}\,dx}$\vspace{1cm}\vfill
\end{enumerate}
\emph{Evaluate the derivative}
\begin{enumerate}[resume]
\item $\frac{d}{dx}$$\frac{2}{3}e^{3x}$\vfill
\item $\frac{d}{dx}\left(\ln(4x+1)+5\right)$\vfill
\end{enumerate}

\end{document}
